% ARCHIVED DRAFT: independently audited by root and incorporated on 2026-11-19 UTC.
% The assembled paper is authoritative; this file retains drafting provenance.
\begin{saresult}{B.076}{Proved}{Principal-conductor pullbacks in arbitrary valuation rings}
Let $V$ be a valuation domain of characteristic $p$ with infinite residue
field $l$, containing a finite subfield $k=\mathbb F_q$. Let
$0\ne x\in\mathfrak m_V$, and put
\[
                 R=k+xV,\qquad\mathfrak m=xV.
\]
For every unital commutative inclusion $R\subset S$, the polynomial
coaddition and comultiplication tensors exist uniquely if and only if
\[
                           (R:_S\mathfrak m)=R.           \tag{B76.1}
\]
This is equivalent to tensor injectivity, to the absence of a
$\mathfrak m$-annihilator in $S/R$, and to $\Omega(R,S)=R[X]$.
All positive cases have every multivariable tensor isomorphism.
No discreteness, rank, completeness, or coefficient-torsion condition
is imposed. The conductor $xV$ is principal as a $V$-ideal; it need
not be principal as an $R$-ideal.
\end{saresult}
\begin{proof}
Let $v$ be the valuation, $\eta=v(x)>0$, and $K=\operatorname{Frac}(V)
=\operatorname{Frac}(R)$. The ring $R$ is local with maximal ideal
$\mathfrak m$ and residue field $k$. The determinant-ideal proof in
B.075 works verbatim with $x$ in place of $t^a$: a grid $xv_i$ with
distinct residues of $v_i\in V$ gives a determinant $x^{d(d+1)/2}$ times
a unit of $V$, and multiplication by $\mathfrak m$ puts a power of
$\mathfrak m$ in its principal $R$-ideal. The opposite cofinality follows
by residue multiplicities. Thus the actual coefficient ring is
\[
                 C=\bigcup_{n\ge0}(R:_S\mathfrak m^n).    \tag{B76.2}
\]
Under tensor injectivity $R\to C$ is epic. Rationalization has image
$L$ satisfying $R\subset L\subset R[1/x]=V[1/x]$, and kernel whose
elements are annihilated by powers of $\mathfrak m$.

There are only two possibilities for this epimorphic overring $L$.
If every $c\in L$ has $v(c)\ge-\eta$, then
$\mathfrak m^2L\subset\mathfrak m\subset R$.
The conductor $J=(R:L)$, an ideal of both rings, therefore contains
$\mathfrak m^2$. The inclusion $R/J\subset L/J$ is epic and its base
has nilpotent maximal ideal and residue field $k$. Reducing modulo
that ideal gives a field epimorphism, hence an isomorphism; finite
iteration of the nilpotent ideal then gives $L/J=R/J$ and $L=R$.
Consequently if $L\ne R$, some $c\in L$ has $v(c)<-\eta$.
Then $c^{-1}\in xV\subset R$ and $c^{-1}V\subset xV$, so $V\subset L$.
Also $x^{-1}=c(xc)^{-1}$ and $(xc)^{-1}\in V$, yielding $L=V[1/x]$.
If $L=R$, rationalization retracts $C$ to $R$ and epicness gives $C=R$.

If $L=V[1/x]$, the polynomial idempotent argument of B.075 applies
unchanged: lift $x^{-1}$ to $c\in C$, use
$x^N(xc-1)=0$ to make $(xc)^N$ idempotent, and split its two factors.
The $x$-invertible factor is $L$, and the other has a power of
$\mathfrak m$ acting by zero. A ring epimorphism from a local ring
with nilpotent maximal ideal is surjective, by the field quotient
and finite iteration just used. Hence
\[
                    C\simeq L\times R/I,
                    \mathfrak m^M\subset I               \tag{B76.3}
\]
for some $M$, allowing $I=R$.

Put $P(X)=X^q-X$ and $U=P(X)/x$. One has
\[
 B=\{F\in K[X]:F(R)\subset V\}
   =V[X,U]/(X^q-X-xU)
   =\bigoplus_{0\le i<q,\ j\ge0}VX^iU^j,
 \quad B/\mathfrak m_VB\simeq\prod_{b\in k}l[U].         \tag{B76.4}
\]
Every polynomial has a finite expansion in this basis over $K$.
If its minimum coefficient valuation were negative, divide by a
coefficient attaining that minimum. The resulting polynomial has
nonzero residue, but its values on $R$ have positive valuation.
Evaluation at $b+xz$, with $z\in V$ varying, has residues
$X=b$, $U=-\bar z$; the infinite field $l$ makes this impossible.
This proves (B76.4), including non-Archimedean value groups.
In particular all these integer-valued polynomials have their
coefficients in $L$.

In this basis the $R$-valued module is
\[
 \Omega(R,K)=
 \bigoplus_{0\le i<q}RX^i
 \ \oplus\!!\bigoplus_{0\le i<q,\ j\ge1}xVX^iU^j.       \tag{B76.5}
\]
To see this directly, subtract the degree-less-than-$q$ interpolation
polynomial with the same values on $k$; it has coefficients in $R$.
The remainder has no $j=0$ coefficients. If its minimum valuation
$\gamma$ is positive but less than $\eta$, all its values lie in
$R\cap\mathfrak m_V=xV$; division by a minimum coefficient then gives
a nonzero residue polynomial vanishing on all $k\times l$, impossible.
If $\gamma=0$, the residue polynomial in each component has values
in the finite field $k$, hence is constant; its zero constant term
makes it zero, again impossible. Thus $\gamma\ge\eta$.
The converse is immediate. If $I=R$ in (B76.3), formula (B76.5) gives
exactly the nonzero tensor kernel $\alpha P\otimes P-P\otimes\alpha P$
for $\alpha\in V$ whose residue lies outside $k$, as in B.075.

For a proper $I$ in (B76.3), let
$H=\{F\in L[X]:F(R)\subset I\}$ and $E=R[X]+H$, the first projection of
$\Omega(R,L\times R/I)$. We prove $H\not\subset\mathfrak mE$ by an
integer-bound argument which requires no least valuation in an ideal.
For $F\in H$, subtract its degree-less-than-$q$ interpolation on $k$;
that polynomial has coefficients in $I$. Denote the resulting module
of polynomials with all $j=0$ coefficients zero by $H_0$.
Every nonzero member of $H_0$ has minimum coefficient valuation at
least $\eta$, by (B76.5). Moreover $x^MU\in H_0$ by (B76.3).
Choose the least positive integer $n$ for which some nonzero
$G\in H_0$ has minimum coefficient valuation $\gamma\le n\eta$.
This integer exists, and it is a minimum in the ordinary integers,
not an assumed minimum in the valuation group.

After division by a minimum coefficient, $G$ has a nonzero residue
component $h(U)$ at some $b_0\in k$. Since its $j=0$ coefficients
vanish, $h(0)=0$ and $h$ is nonconstant.
Precompose with $\psi_N(X)=b_0+xU(X)^N$, for an integer $N>n$.
This maps $R$ into $R$ and preserves vanishing on $k$, so
$F=G\circ\psi_N\in H_0$.
The minimum coefficient valuation of $F$ is still $\gamma$, and its
normalized residue is $h(-U^N)$ in every component. Hence some
coefficient of $U$-degree at least $N$ has valuation $\gamma$.

If $H\subset\mathfrak mE$, write
$F=A+\sum_i r_iF_i$, where $A\in\mathfrak mR[X]$,
$r_i\in\mathfrak m$, and $F_i\in H$.
Subtract the interpolation polynomials of the $F_i$ and absorb their
products into $A$, so that every $F_i\in H_0$.
The coefficient of $X^iU^j$ in $\mathfrak mR[X]$ lies in
$x^{j+1}V$, since the corresponding monic $R$-basis is $X^iP^j$.
At the chosen degree $j\ge N>n$ its valuation is therefore greater
than $\gamma$. Some summand $r_iF_i$ must have a coefficient there
of valuation at most $\gamma$. The minimum coefficient valuation
of that nonzero $F_i$ is at most
\[
                  \gamma-v(r_i)\le\gamma-\eta
                                     \le(n-1)\eta,
\]
If $n=1$, this contradicts the lower bound $\eta>0$ already proved
for every nonzero element of $H_0$. If $n>1$, it contradicts the
minimal choice of $n$. Thus $H\not\subset\mathfrak mE$.

Choose $F\in H\setminus\mathfrak mE$. As in B.075,
$f=(F,0)$ survives modulo $\mathfrak m\Omega$ by its first projection.
The polynomial $g=(0,P^M\bmod I)$ is null on $R$ and survives there
by its monic second projection over $k$.
Their pure tensor is nonzero modulo $\mathfrak m$ but its polynomial
image is zero. This excludes every proper $I$ in (B76.3).
We have proved that tensor injectivity forces $C=R$.
The iterated colon identity makes $C=R$ equivalent to (B76.1), and
B.064 proves all remaining positive assertions.
\end{proof}
